\section{Mixed centered harmonic-tail moments}
\label{mt:section}

The source reports evaluate every centered even moment of a trigamma
square \cite{IncomingHigherReflection,IncomingParity}; the harmonic-parity
report asks for the corresponding evaluation for two unequal even-order
harmonic tails. We resolve that two-factor problem for all
orders and all moments. The resulting finite formula reduces entirely to
ordinary zeta values and their products. We then package the infinite
moment sequence into one analytic generator with explicit
polylogarithmic primitives. The centered parity mechanism is inherited
from the source criterion; the evaluation and the generator are the
additional results here. Euler's classical odd-weight reduction supplies
the arithmetic step~\cite{Bradley2007}.

\subsection{Continuation and ordinary convergent representatives}

Let $a,b\ge2$ be even integers, and put
\begin{equation}\label{mt:notation}
 T_a(n)=\zeta(a,n+1)=\zeta(a)-H_n^{(a)}
       =\frac{\psi^{(a-1)}(n+1)}{(a-1)!},
 \qquad x_n=n+\frac12.
\end{equation}
The last equality uses that $a$ is even. Define
\begin{equation}\label{mt:F}
 F_{a,b}(s)=\sum_{n=0}^{\infty}\frac{T_a(n)T_b(n)}{x_n^s},
 \qquad V_{a,b}(M)=F_{a,b}(-2M),\quad M\ge0.
\end{equation}
The series for $F_{a,b}$ converges absolutely when
$\Re s>3-a-b$. The value at $-2M$ is understood through the continuation
proved below whenever the unmodified series diverges.

Write $\alpha=a+b-2$, and define rational coefficients
\begin{align}
 t_{a,0}&=\frac1{a-1},&
 t_{a,k}&=\frac{B_{2k}(1/2)}{(2k)!}(a)_{2k-1}\quad(k\ge1),
 &c_k&=\sum_{j=0}^k t_{a,j}t_{b,k-j}.
 \label{mt:centered-coefficients}
\end{align}
Here $(a)_r$ is the rising factorial, $B_j(x)$ is the Bernoulli
polynomial, and $B_j=B_j(0)$, with $B_1=-1/2$.
The centered Hurwitz expansions \cite[\S25.11(xii)]{DLMFHurwitz} give
\begin{equation}\label{mt:centered-expansion}
 T_a(n)\sim\sum_{k\ge0}t_{a,k}x_n^{1-a-2k},
 \qquad
 T_a(n)T_b(n)\sim\sum_{k\ge0}c_kx_n^{-\alpha-2k}.
\end{equation}
Only finite truncations of these asymptotic expansions are used.

\begin{proposition}[Regularity and bracketed moments]
\label{mt:regularity}
The function $F_{a,b}$ extends meromorphically to the complex plane.
Its only possible poles are simple poles at
$s=3-a-b-2k$, $k\ge0$, with residue $c_k$.
In particular every even spectral integer is regular. For
\begin{equation}\label{mt:subtraction-polynomial}
 P_M(x)=\sum_{\alpha+2k\le2M}c_kx^{2M-\alpha-2k},
\end{equation}
one has the ordinary convergent identity
\begin{equation}\label{mt:bracket}
 V_{a,b}(M)=\sum_{n=0}^{\infty}
 \left[x_n^{2M}T_a(n)T_b(n)-P_M(x_n)\right].
\end{equation}
The polynomial is zero if its indexing set is empty. The bracket is
$O(x_n^{-2})$.
\end{proposition}

\begin{proof}
For any fixed $K$, subtract the first $K+1$ terms of
\eqref{mt:centered-expansion} inside the defining series. The remainder
converges normally in a right half-plane extending farther to the left,
while the removed terms
sum to
$\sum_{k=0}^K c_k\zeta(s+\alpha+2k,1/2)$.
Increasing $K$ provides compatible continuations and the stated poles
and residues. Since $a+b$ is even, all these potential poles are odd
integers. At $s=-2M$, every nondecaying term has a nonnegative even
power of $x_n$. Its continued sum is
$\zeta(-2j,1/2)=0$, including $j=0$. After their removal the next
possible power is $x_n^{-2}$, which proves \eqref{mt:bracket}.
\end{proof}

\begin{corollary}[Every spectral derivative with its correction terms]
\label{mt:spectralderivatives}
For $r,M\geq0$, every derivative at a centered even spectral value has
the ordinary convergent representation
\begin{align}
 F_{a,b}^{(r)}(-2M)
  ={}&\sum_{n=0}^{\infty}(-\log x_n)^r
       \left[x_n^{2M}T_a(n)T_b(n)-P_M(x_n)\right]\notag\\
    &+\sum_{\alpha+2k\leq2M}
               c_k\zeta^{(r)}(\alpha+2k-2M,1/2).
 \label{mt:logmoment}
\end{align}
The second sum vanishes for $r=0$. It must in general be retained for
$r\geq1$.
\end{corollary}
\begin{proof}
Use the same finite subtractions as in \eqref{mt:subtraction-polynomial}
at a moving $s$ in a neighborhood of $-2M$. The remainder series and all
its fixed spectral derivatives converge locally normally there, since
their terms are bounded by $C x_n^{-2+\varepsilon}(1+|\log x_n|^r)$
with some $\varepsilon<1$. Differentiation of the finite Hurwitz
correction gives exactly \eqref{mt:logmoment}.
\end{proof}

For example, the first derivative in the $(2,4)$ family at $s=-4$ is
\begin{equation}\label{mt:derivative-example}
 F'_{2,4}(-4)=
 -\sum_{n\geq0}\log x_n
        \left[x_n^4T_2(n)T_4(n)-\frac13\right]
 -\frac{\log 2}{6},
\end{equation}
because $\zeta'(0,1/2)=-\tfrac12\log2$. This identity fixes the
logarithmic remainder and its constant; it does not claim that the
remainder is an ordinary-zeta combination.

We will also use a concrete equivalent characterization. Let
$\operatorname{CT}_N$ mean the coefficient of $N^0(\log N)^0$ in
an asymptotic expansion in the integer cutoff $N$. Then
\begin{equation}\label{mt:cutoff-value}
 V_{a,b}(M)=\operatorname{CT}_N
   \sum_{n=0}^{N-1}x_n^{2M}T_a(n)T_b(n).
\end{equation}
Indeed the partial sum of each removed even polynomial power is
$B_{2j+1}(N+1/2)/(2j+1)$, an odd polynomial in $N$ with zero constant
term. The summable remainder in \eqref{mt:bracket} tends to its ordinary
sum. The coordinate $N$ in \eqref{mt:cutoff-value} is part of the
statement; no change of endpoint coordinate is implicit.

\subsection{A finite formula at every moment order}

The double-zeta convention used below is
\[
 \zeta(u,v)=\sum_{j>k\ge1}\frac1{j^u k^v}.
\]
Thus the first exponent belongs to the larger summation variable.
For each positive odd integer $d$, define
\begin{equation}\label{mt:oriented}
 K_{a,b}(d)=
 \begin{cases}
  \displaystyle\zeta(b,a-d)+\frac12\zeta(a+b-d),&d<a,\\[2mm]
  \displaystyle\sum_{j=1}^{\min(b/2,(d-a+1)/2)}
    \frac{\binom{d-a+1}{2j}B_{d-a+1-2j}}{d-a+1}
             \zeta(b-2j),&d>a,
 \end{cases}
\end{equation}
and put
\begin{equation}\label{mt:E}
 E_{a,b}(d)=K_{a,b}(d)+K_{b,a}(d).
\end{equation}
There is no missing case $d=a$: $d$ is odd and $a$ even. All double
zetas in the first branch are ordinary convergent sums. In the second
branch, $\zeta(0)=-1/2$ has its usual continued meaning.

\begin{theorem}[All mixed even-tail moments]
\label{mt:all-moments}
For even $a,b\ge2$ and every integer $M\ge0$,
\begin{equation}\label{mt:moment-formula}
 \boxed{\displaystyle
 V_{a,b}(M)=\frac1{2M+1}\sum_{\ell=0}^{M}
  \binom{2M+1}{2\ell+1}B_{2M-2\ell}(1/2)
                          E_{a,b}(2\ell+1).}
\end{equation}
The summation and every coefficient in this formula are finite and
explicit.
\end{theorem}

The $j=b/2$ term in \eqref{mt:oriented}, whenever its range includes
it, records an exact rational boundary contribution through
$\zeta(0)=-1/2$. Discarding this term would already give an incorrect
answer for divergent centered moments of the trigamma square.

\begin{proof}
First form the exact polynomial partial sum
\begin{align}
 S_M(k)&=\sum_{n=0}^{k-1}(n+1/2)^{2M}
       =\frac{B_{2M+1}(k+1/2)}{2M+1}
       =\sum_{\ell=0}^{M}s_{M,\ell}k^{2\ell+1},
 \label{mt:S}\\
 s_{M,\ell}&=\frac{\binom{2M+1}{2\ell+1}}{2M+1}
                         B_{2M-2\ell}(1/2).
 \notag
\end{align}
The tail recurrence gives
\begin{align*}
&T_a(k-1)T_b(k-1)-T_a(k)T_b(k)\\
&\hspace{18mm}=k^{-a}T_b(k-1)+k^{-b}T_a(k-1)-k^{-a-b}.
\end{align*}
Applying finite summation by parts, without taking any divergent limits,
therefore yields
\begin{equation}\label{mt:finite-sbp}
 \sum_{n=0}^{N-1}x_n^{2M}T_a(n)T_b(n)
 =\sum_{\ell=0}^{M}s_{M,\ell}\,\mathcal E_{2\ell+1}(N),
\end{equation}
where
\begin{align}
 \mathcal E_d(N)={}&N^dT_a(N)T_b(N)
   +\sum_{k=1}^N k^{d-a}T_b(k-1)
   +\sum_{k=1}^N k^{d-b}T_a(k-1)
   -H_N^{(a+b-d)}.
 \label{mt:finite-E}
\end{align}
We compute its cutoff constant explicitly.

If $d<a$, the first oriented sum in \eqref{mt:finite-E} is absolutely
convergent, and direct rearrangement gives
\begin{equation}\label{mt:convergent-oriented}
 \sum_{k=1}^{\infty}k^{d-a}T_b(k-1)
 =\zeta(b,a-d)+\zeta(a+b-d).
\end{equation}
If $d>a$, put $p=d-a$, a positive odd integer. The Faulhaber
polynomial is
\begin{equation}\label{mt:faulhaber}
 Q_p(X)=\sum_{k=1}^X k^p
   =\frac12X^p+\sum_{j=1}^{(p+1)/2}q_{p,2j}X^{2j},
 \qquad
 q_{p,2j}=\frac{\binom{p+1}{2j}B_{p+1-2j}}{p+1}.
\end{equation}
The finite-sum notation specifies the polynomial by its values at
nonnegative integers. Define $q_{p,h}=0$ outside its displayed
nonzero even degrees, whenever an even $h$ is used below. A finite
rearrangement which retains the infinite tail gives the exact equality
\begin{equation}\label{mt:finite-rearrangement}
 \sum_{k=1}^N k^pT_b(k-1)
  =\sum_{h\ge1}[X^h]Q_p(X)\,H_N^{(b-h)}
    +Q_p(N)T_b(N).
\end{equation}
For $v>1$, the cutoff constant of $H_N^{(v)}$ is $\zeta(v)$;
for $v=1$ it is Euler's constant $\gamma$; and for integer $v\le0$
it is zero because the finite power sum is a polynomial in $N$ with
zero constant term.

At the potential logarithmic collision $d=a+b-1$, the half-leading
monomial in each of the two Faulhaber polynomials contributes
$\gamma/2$. The diagonal term in \eqref{mt:finite-E} contributes
$-\gamma$. These cancel. In particular, this proof never evaluates
an isolated $\zeta(1)$.

We now evaluate the constant contributions of the explicit tails. Write
\begin{equation}\label{mt:tau}
 T_a(N)\sim\sum_j\tau_a(j)N^{-j},\qquad
 \begin{cases}
 \tau_a(a-1)=1/(a-1),\\
 \tau_a(a)=-1/2,\\
 \tau_a(a+2k-1)=B_{2k}(a)_{2k-1}/(2k)!,&k\ge1,
 \end{cases}
\end{equation}
with all other coefficients zero. Since $a,b$ are even and $d$ odd,
an odd total degree in the product can only use one exceptional even
index $a$ or $b$. Hence
\begin{equation}\label{mt:boundary-product}
 \operatorname{CT}_N N^dT_a(N)T_b(N)
 =-\frac12\tau_b(d-a)-\frac12\tau_a(d-b).
\end{equation}
For $p=d-a>0$, the only odd monomial of $Q_p$ is $X^p/2$,
and the only even index in $T_b$ is $b$. Therefore
\begin{equation}\label{mt:boundary-oriented}
 \operatorname{CT}_N Q_p(N)T_b(N)
 =\frac12\tau_b(p)-\frac12q_{p,b}.
\end{equation}
The $\tau$ terms in \eqref{mt:boundary-product} cancel those from
the two instances of \eqref{mt:boundary-oriented}. The remaining
constant is
\begin{equation}\label{mt:boundary-constant}
 -\frac12\bigl(q_{d-a,b}+q_{d-b,a}\bigr),
\end{equation}
where a term is zero if its first index is not positive. This is
exactly the contribution of the $\zeta(0)$ terms in
\eqref{mt:oriented}.

For completeness consider the remaining diagonal constant. If
$W=a+b-d\ge3$, a branch $d<a$ contributes the full $\zeta(W)$
in \eqref{mt:convergent-oriented}, and a branch $d>a$ contributes
$\zeta(W)/2$ from the half-leading term in
\eqref{mt:finite-rearrangement}. After subtracting the diagonal
$\zeta(W)$, exactly one half remains for each branch $d<a$.
When $W=1$ the constants cancel as already shown. When $W<1$,
the diagonal power sum has zero constant term. The other surviving
Faulhaber terms have even degree and yield the even zeta values in
\eqref{mt:oriented}. Thus
\[
 \operatorname{CT}_N\mathcal E_d(N)=E_{a,b}(d).
\]
Taking constants in \eqref{mt:finite-sbp} and applying
\eqref{mt:cutoff-value} proves \eqref{mt:moment-formula}.
\end{proof}

\subsection{Reduction to ordinary zeta values}

For clarity, the classical reduction needed to eliminate every double
coordinate is written out. Let $m\ge2$ be even, let $n\ge1$ be odd,
and put $w=m+n$ and
$\mathcal U_{m,n}=\zeta(m,n)+\zeta(w)/2$. Euler's formulas, in the
convention above, give
\begin{equation}\label{mt:euler-one}
 \mathcal U_{m,1}=\frac{m+1}{2}\zeta(m+1)
  -\sum_{j=1}^{m/2-1}\zeta(2j)\zeta(m+1-2j),
\end{equation}
and, for $n\ge3$,
\begin{align}
 \mathcal U_{m,n}={}&\frac12\binom wm\zeta(w)+\zeta(m)\zeta(n)
 \notag\\
 &-\sum_{j=1}^{(w-3)/2}
  \left[\binom{w-2j-1}{m-1}+\binom{w-2j-1}{n-1}\right]
                      \zeta(2j)\zeta(w-2j).
 \label{mt:euler-odd}
\end{align}
These are the ordinary specialization of the classical signed Euler
reduction; an elementary proof from partial fractions is provided in
\cite[Eq.~(2)--(3) and Proposition~1]{Bradley2007}. No divergent
$\zeta(1)$ appears in \eqref{mt:euler-one}--\eqref{mt:euler-odd}.
The first branch of \eqref{mt:oriented} is exactly
$\mathcal U_{b,a-d}$, so the formulas apply to every such branch.

\begin{corollary}[Ordinary-zeta closure and a fixed spanning list]
\label{mt:closure}
Every $V_{a,b}(M)$ is a rational linear combination of $1$, ordinary
even zeta values, ordinary odd zeta values, and products of one even
and one odd zeta value. Put $B=\max(a,b)$. The entire moment sequence
lies in the rational span of the at most $B$ coordinates
\begin{equation}\label{mt:finite-span}
 E_{a,b}(1),E_{a,b}(3),\ldots,E_{a,b}(B-1),
 \qquad 1,\zeta(2),\zeta(4),\ldots,\zeta(B-2).
\end{equation}
The positive even-zeta sublist is empty when $B=2$.
\end{corollary}
\begin{proof}
Ordinary-zeta closure follows by inserting
\eqref{mt:euler-one}--\eqref{mt:euler-odd} into
\eqref{mt:oriented}. For odd $d>B$, both branches in
\eqref{mt:oriented} are finite combinations of
$\zeta(0),\zeta(2),\ldots,\zeta(B-2)$. Formula
\eqref{mt:moment-formula} then proves the spanning assertion.
\end{proof}

The corollary is a finite spanning theorem, not a claim that its
coordinates are arithmetically independent or minimal. It answers the
ordinary-constant reduction question for every pair of even-order
tails; increasing the moment order introduces no additional constants.

\subsection{The first unequal-order family}

For $a=2$ and $b=4$, the first five coordinates are
\begin{align}
 E_{2,4}(1)&=\frac{15}{2}\zeta(5)-3\zeta(2)\zeta(3),
 &E_{2,4}(3)&=\frac12\zeta(2)+\frac32\zeta(3),\notag\\
 E_{2,4}(5)&=\frac14\zeta(2)-\frac38,
 &E_{2,4}(7)&=-\frac1{12}\zeta(2)-\frac13,\label{mt:E24}\\
 E_{2,4}(9)&=\frac1{12}\zeta(2)+\frac3{16}.\notag
\end{align}
The centered product expansion begins
\begin{equation}\label{mt:T24-expansion}
 T_2(n)T_4(n)=\frac1{3x_n^4}-\frac7{36x_n^6}
                       +\frac{61}{360x_n^8}+O(x_n^{-10}).
\end{equation}
Theorem~\ref{mt:all-moments} consequently gives the following ordinary
convergent identities:
\begin{align}
 \sum_{n\ge0}T_2(n)T_4(n)
  &=\frac{15}{2}\zeta(5)-3\zeta(2)\zeta(3),
 \label{mt:V0}\\
 \sum_{n\ge0}x_n^2T_2(n)T_4(n)
  &=\frac14\zeta(2)\zeta(3)+\frac16\zeta(2)+\frac12\zeta(3)
                   -\frac58\zeta(5),
 \label{mt:V1}\\
 \sum_{n\ge0}\left[x_n^4T_2(n)T_4(n)-\frac13\right]
  &=-\frac7{80}\zeta(2)\zeta(3)-\frac1{30}\zeta(2)
      -\frac14\zeta(3)+\frac7{32}\zeta(5)-\frac3{40}.
 \label{mt:V2}
\end{align}
One further example is
\begin{align}
 &\sum_{n\ge0}\left[x_n^6T_2(n)T_4(n)-\frac{x_n^2}{3}+\frac7{36}\right]
 \notag\\
 &\qquad=\frac{31}{448}\zeta(2)\zeta(3)-\frac1{672}\zeta(2)
          +\frac7{32}\zeta(3)-\frac{155}{896}\zeta(5)
          +\frac{31}{672}.
 \label{mt:V3}
\end{align}
Since $T_2=\psi'$ and $T_4=\psi'''/6$, multiplication by $6$
turns these into identities for products of two unequal polygamma
orders. They equally evaluate the centered harmonic numerator
$(H_n^{(2)}-\zeta(2))(H_n^{(4)}-\zeta(4))$.

\subsection{An analytic generator for all moment orders}

For a fixed ordered pair $(a,b)$, define
\begin{equation}\label{mt:G}
 G_b(t)=\frac12\coth(t/2)
             \sum_{j=1}^{b/2}\zeta(b-2j)\frac{t^{2j}}{(2j)!}.
\end{equation}
Its apparent singularity at zero is removable. It is odd and analytic
on $|t|<2\pi$.

\begin{theorem}[All-moment generating function]
\label{mt:generator}
Let $A_{a,b}$ be the unique analytic solution on $|t|<2\pi$ of
\begin{equation}\label{mt:generator-ode}
 A_{a,b}^{(a)}(t)=G_b(t),
\end{equation}
whose Taylor polynomial of degree less than $a$ is
\begin{equation}\label{mt:initial-data}
 \sum_{\substack{1\le d<a\\d\text{ odd}}}
  \left[\zeta(b,a-d)+\frac12\zeta(a+b-d)\right]\frac{t^d}{d!}.
\end{equation}
Then the exponential series
\begin{equation}\label{mt:moment-egf}
 \mathscr V_{a,b}(t)=\sum_{M=0}^{\infty}
                      V_{a,b}(M)\frac{t^{2M}}{(2M)!}
\end{equation}
converges for $|t|<2\pi$ and satisfies
\begin{equation}\label{mt:generator-formula}
 \boxed{\displaystyle
 \mathscr V_{a,b}(t)=
       \frac{A_{a,b}(t)+A_{b,a}(t)}{2\sinh(t/2)}.}
\end{equation}
The quotient has a removable value at zero and is even analytic.
\end{theorem}

\begin{proof}
Anchored integration of the analytic function $G_b$ gives existence
and uniqueness. Its oddness, and the even value of $a$, show that
$A_{a,b}$ is odd. The Bernoulli generating identity
\[
 C(t)=\frac t2\coth(t/2)
       =\sum_{k\ge0}\frac{B_{2k}}{(2k)!}t^{2k}
\]
gives
\[
 G_b(t)=C(t)\sum_{j=1}^{b/2}\zeta(b-2j)\frac{t^{2j-1}}{(2j)!}.
\]
For positive odd $p$, its coefficient of $t^p/p!$ is
\[
 \sum_{j=1}^{\min(b/2,(p+1)/2)}
   \frac{\binom{p+1}{2j}B_{p+1-2j}}{p+1}\zeta(b-2j).
\]
Together with the initial data, this says exactly that
\begin{equation}\label{mt:A-series}
 A_{a,b}(t)=\sum_{\substack{d\ge1\\d\text{ odd}}}
                             K_{a,b}(d)\frac{t^d}{d!}.
\end{equation}
Finally,
\[
 \frac1{2\sinh(t/2)}
     =\frac1t\sum_{k\ge0}\frac{B_{2k}(1/2)}{(2k)!}t^{2k}.
\]
Multiplication with the sum of the two series in
\eqref{mt:A-series} gives precisely the coefficient formula
\eqref{mt:moment-formula}. The numerator is odd and vanishes at zero;
the denominator has a simple zero there and no other zero in this
disk. Hence the quotient is even analytic throughout the disk, proving
actual convergence of \eqref{mt:moment-egf} as well as the identity.
\end{proof}

This analytic generator concerns the regularized values already fixed
by \eqref{mt:bracket}. It is not obtained by interchanging an exponential
series with the divergent unmodified moment sums.

\subsection{Explicit polylogarithmic primitives and their constants}

For an integer $k\ge1$, set
\begin{equation}\label{mt:L-integral}
 L_k(t)=\int_0^t\frac{u^k}{e^u-1}\,du.
\end{equation}
The integrand is analytic at zero. Initially in $\Re t>0$, the standard
polylogarithm differentiation rule \cite[\S25.12]{DLMFPolylog} gives the explicit anchored formula
\begin{equation}\label{mt:L-polylog}
 L_k(t)=k!\zeta(k+1)
 -\sum_{h=0}^k\frac{k!}{(k-h)!}
                      t^{k-h}\operatorname{Li}_{h+1}(e^{-t}).
\end{equation}
For real $t>0$, all polylogarithms have arguments in $(0,1)$; this fixes
the branch. The formula extends to the right half-plane by analyticity.
Differentiating the finite sum makes its adjacent terms cancel and
leaves $t^k/(e^t-1)$. As $t\to0$ from the right, all positive powers
multiplying the polylogarithms vanish, including
$t^k\operatorname{Li}_1(e^{-t})$, and the remaining term cancels
$k!\zeta(k+1)$. This proves both the derivative and the anchoring
constant. The completed expression, taken as a whole, then continues
analytically to $|t|<2\pi$ through its integral definition. Individual
polylogarithm terms need not be assigned separate analytic values
around the logarithmic point $t=0$.

For positive even $m$, define
\begin{align}
 R_{a,m}(t)={}&\frac{t^{a+m}}{2(a+m)!}
 \notag\\
 &+\frac1{m!(a-1)!}\sum_{r=0}^{a-1}(-1)^r\binom{a-1}{r}
                                  t^{a-1-r}L_{m+r}(t).
 \label{mt:R}
\end{align}
Expanding the anchored integration kernel $(t-u)^{a-1}$ shows that
\begin{equation}\label{mt:R-integral}
 R_{a,m}(t)=\frac1{(a-1)!}\int_0^t(t-u)^{a-1}
                         \frac12\coth(u/2)\frac{u^m}{m!}\,du.
\end{equation}
Here one uses $\tfrac12\coth(u/2)=\tfrac12+(e^u-1)^{-1}$;
the first term gives the polynomial in \eqref{mt:R}, and the second
gives the finite sum of $L$ integrals. Consequently
\begin{equation}\label{mt:A-polylog}
 A_{a,b}(t)=(\text{the polynomial in \eqref{mt:initial-data}})
              +\sum_{j=1}^{b/2}\zeta(b-2j)R_{a,2j}(t).
\end{equation}
Equations \eqref{mt:euler-one}--\eqref{mt:euler-odd} and
\eqref{mt:generator-formula}--\eqref{mt:A-polylog} express the whole
moment sequence using ordinary zeta constants, elementary functions,
and finitely many $\operatorname{Li}_k(e^{-t})$ with $k\le a+b$.
Every repeated primitive is anchored at zero. Near zero, the ODE or
Taylor form is more suitable for numerical evaluation because the
separate polylogarithmic terms undergo strong cancellation.

\subsection{Scope of the result}

The theorem resolves the two-factor even-tail sector of the source
question. It does not assert ordinary-zeta closure for the first
spectral derivative $F'_{a,b}(-2M)$, independent order derivatives,
or products of four tails. In particular, adding a spectral derivative
introduces a global logarithmic moment; the finite-value theorem alone
does not eliminate it. Verification records and precise next questions
are given in Sections~\ref{audit:section} and~\ref{research:section}.
